\chapter{Main Theorem}

Now that we have presented and proved the existence of fractional Gaussian fields $X_s$ on $K$ and discrete fractional Gaussian fields $X_s^{m}$ on $V_m$ out next task is proving they converge. Therefore in this section our main goal is to show that $X_s^{m}$ converges in distribution to $X_s$.

\section{Preliminary lemmas}
In this section we will prove a collection of lemmas that will be needed later.

\begin{lemma}\label{suckItKigami}
	Let $\cbr{P_t}_{t \geq 0}$ be the semigroup generated by the Dirichlet Laplacian $(-\Delta)^{-s}$ and $p_t(x,y)$ the heat kernel admitted by $\cbr{P_t}_{t \geq 0}$. Then if $\phi_i$ is an orthonormal eigenfunction of $(-\Delta)^{-s}$ then
	\begin{equation}
		P_t \phi_i(x) = e^{-\lambda_i t}\phi_i(x)
	\end{equation}
\end{lemma}
\begin{proof}
	 It follows immediately from the calculations:
	\begin{multline}
		P_t \phi_i(x) = \int_{K}p_t(x,y)\phi_i(y) \dd \mu(y) = \int_{K} \left(\sum_{j = 1}^{\infty}e^{-\lambda_j} \phi_j(x)\phi_j(y)\right) \phi_i(y) \dd \mu(y) \\
		= \int_{K} e^{-\lambda_i t}\phi_i(x) \dd \mu(y) = e^{-\lambda_i t}\phi_i(x)
	\end{multline}
\end{proof}

In the following proof we will use results from \cite{KURTZ1969354}. Before stating the lemma and proof we will give some arguments as to why we are allowed to use these results.

Firstly we note that since $\Delta$ is a bounded operator the by \cite[Theorem 1.2]{Pazy} it is the infinitesimal generator of a uniformly continuos semigroup.

Secondly we have that the sequence of \textcolor{red}{laplacians} $\cbr{\Delta_m}_{m \geq 0}$ indeed coincides with extended limit defined in \cite[page 355]{KURTZ1969354}. This follows from \cref{pointwiseFormula}.

Lastly the range $\mathcal{R}(\lambda_0 - \Delta)$ is dense in $L^2$ \textcolor{red}{why?}


\begin{lemma}\label{convergenceOfSemigroups}
	For all $f \in S(K)$ and $t > 0$,
	\begin{equation}
		\lim_{m \to \infty} \frac{1}{a_m} \sum_{p \in V_m} f_m(p)P_t^{m} f_m(p) = \int_{K}f(x)P_tf(x) \dd\mu (x)
	\end{equation}
	where $f_m = f \mid_{V_m}$.
\end{lemma}
\begin{proof}
	In this proof we wish to use \cite[theorem 2.1]{KURTZ1969354}  so recall that by the definition of the Laplacian on $K$ we have that
	\begin{equation}
		\lim_{m \to \infty} \sup_{p \in V_m} \abs{\Delta_{m}f_m(p) -\Delta f(p)}.
	\end{equation}
	Thus by \cite[theorem 2.1]{KURTZ1969354} \textcolor{red}{Is defined with semigroup T(s) and page 354, is $P_n$ in our case $\int_{K}\psi_{x}^{(m)} \dd \mu$} we have that for all $t \geq 0$
	\begin{equation}\label{convergensofPt}
		\lim_{m \to \infty} \sup_{p \in V_m} \abs{P_t^{m}f_m(p) - P_t f(p)} = 0.
	\end{equation}
	So if we now write 
	\begin{equation}
		\frac{1}{a_m} \sum_{p \in V_m} f_m(p) P_t^{m}f_m(p) = \int_{V_m} f_m P_t^{m} f_m \dd \mu_m.
	\end{equation}
	Thus 
	\begin{equation}
		\int_{V_m} f_m P_t^{m} f_m \dd\mu_m = 		\int_{V_m} f_m \left(P_t^{m} f_m - (P_t f)_m \right) \dd\mu_m + \int_{V_m} f_m (P_t f)_m \dd \mu_m .
	\end{equation}
	The first integral converges to $0$ as $m\to \infty$ because of \eqref{convergensofPt}. Further by lemma we have \newline $\int_{V_m} f_m (P_t f)_m \dd \mu_m \to \int_{K} fP_tf \dd \mu$ as $m \to \infty$.
\end{proof}

\begin{lemma}\label{convergence}
	For all $f \in S(K)$ and $s \geq 0$ the following holds
	\begin{equation}
		\frac{1}{a_m} \sum_{p \in V_m} f_m(p) (-\Delta_m)^{-2s}f_m(p) \to \int_{K}f(x)(-\Delta)^{-2s}f(x) \dd \mu(x), \quad \text{as } m \to \infty
	\end{equation}
\end{lemma}
\begin{proof}
	For $s = 0$  we find that 
	\begin{align}
		(-\Delta_m)^{-0} f&= \sum_{i = 1}^{N_m} \frac{1}{(\lambda_i^{m})^{0}} \phi_{i}^{m}(x) \frac{1}{a_m} \sum_{y \in V_m} \psi_{i}^{m}(y) f(y) \\
		&= \sum_{i = 1}^{N_m} \phi_{i}^{m}(x) \int_{V_m} \phi_i^{m}(y)f(y) \dd \mu_{m} \\
		& = f_m.
	\end{align}
	Likewise by definition of $(-\Delta)^{-s}$ we find
	\begin{align}
		(-\Delta)^{-0} f = \sum_{i = 1}^{\infty} \phi_{i} \int_{K} \phi_i(y) f(y) \dd \mu = f.
	\end{align}
	and the lemma follows from \cref{convergenceOfMeasures}. Now assume $s > 0$ then by Fubini's theorem and teh fact that the semi-group is linear we have that
	\begin{align}
		\frac{1}{a_m} \sum_{p \in V_m} f_m(p) P_t^{m} f_m(p) (-\Delta_{m})^{-2s}f_m(p)  = \frac{1}{\Gamma(2s)} \int_{0}^{\infty} t^{2s-1} \frac{1}{a_m} \sum_{p \in V_m} f_m(p) P_{t}^{m} f_m(p).
	\end{align}
	Now by \cref{suckItKigami} and the definition of $f_m$ we find that
	\begin{align}
		\frac{1}{a_m} \sum_{p \in V_m} f_m(p) P_{t}^{m} f_m(p) &= \frac{1}{a_m} \sum_{p \in V_m} f_m(p) P_{t}^{m} \sum_{j =1}^{\infty} \inn{f_m}{\phi_i} \phi_{i} \\
		&= \frac{1}{a_m} \sum_{p \in V_m} f_m(p) \sum_{j =1}^{\infty} \inn{f_m}{\phi_j} P_{t}^{m} \phi_{j} \\
		& = \frac{1}{a_m} \sum_{p \in V_m} f_m(p) \sum_{j =1}^{\infty} \inn{f_m}{\phi_j} e^{-\lambda_{j}^{m} t} \phi_{j} \\
		&\leq e^{-\lambda_1^{m}t} \frac{1}{a_m} \sum_{p \in V_m} f_m(p)^2
	\end{align}
	where we can guarantee the inequality for all $j$ since $0 < \lambda_1 < \lambda_2 < .\lambda_3 <..$. 
	
	Note that $\sup_{m} \frac{1}{a_m} \sum_{p \in V_m} f_m(p)^2 < \infty$ since $f \in S(K) \subset C(K)$ so $f$ is bounded and $\# V_m < \infty$ for all $m$. We can now by dominated convergence, Fubini's theorem and \cref{convergenceOfSemigroups} find that 
	\begin{align}
		\lim_{m \to \infty} \frac{1}{a_m} \sum_{p \in V_m} f_m(p) (-\Delta_{m})^{-2s} f_m(p) &= \frac{1}{\Gamma(2s)} \int_{0}^{\infty} t^{2s-1} t^{2s-1} \lim_{m \to \infty}\frac{1}{a_m} \sum_{p \in V_m} f_m(p) P_{t}^{m} f_m(p) \dd t \\
		&= \frac{1}{\Gamma(2s)} \int_{0}^{\infty} t^{2s-1} \int_{K} f(x) P_t f(x) \dd\mu(x) \dd t \\
		&= \frac{1}{\Gamma(2s)} \int_{K} f(x) \int_{0}^{\infty} t^{2s-1}P_t f(x) \dd t \dd\mu(x) \\
		&= \int_{K} f(x) (-\Delta)^{-2s}f(x) \dd \mu(x)
	\end{align}
	as desired.
\end{proof}


\section{Main Theorem}
With everything in mind we can prove the convergence in distribution of discrete fractional Gaussian fields to fractional Gaussian fields.

\begin{theorem}
	Let $s \geq 0$, then $X_{s}^{m}$ converges to $X_s$ in distribution in $S'(K)$ when $m \to \infty$
\end{theorem}

\begin{proof}
	First since $S(K)$ is a nuclear space, by \cite[Corollary~2.4]{biermé2017generalizedrandomfieldslevys} it is enough to prove the convergence of the characteristic functional, e.i that for all $f \in S(K)$
	\begin{equation}
		\e \left[\exp(itX_{s}^{m}(f_m))\right] \to \e \left[\exp(itX_{s}(f))\right], \quad \text{ as } m \to \infty
	\end{equation} 
	where $f_m = f \mid_{V_m}$. It follows from \cref{theOneWeNeedInTheEndSomehow} that
	\begin{equation}
		\e \left[\exp(itX_{s}^{m}(f_m))\right] = \exp \left(-\frac{1}{2} t^2 \e \left[(X_{s}^{m}(f_m))^2\right]\right) = \exp \left(-\frac{1}{2a_m} t^2  \sum_{p \in V_m} f_m(p) (-\Delta_m)^{-2s} f_m(p)\right)
	\end{equation}
	Likewise we find that via \cref{fracGaussianField} we get that
	\begin{equation}
		\e \left[\exp(i t X_{s}(f))\right] = \exp\left(-\frac{1}{2} t^2 \e \left[(X_{s}(f))^2\right]\right) = \exp\left(-\frac{1}{2}t^2 \int_{K} f (-\Delta)^{-2s}f \dd \mu\right).
	\end{equation} 
	Where we notice that from \cref{coroForRieszKernelAndFractionalLaplacian}
	\begin{align}
		\e \left[(X_{s}(f))^2\right] &= \int_{K} (-\Delta)^{s}f (-\Delta)^{s}f \dd \mu \\
		& = \int_{K} \int_{K} f(p)f(q) G_{2s}(p,q) \dd\mu(p) \dd \mu(q) \\
		& = \int_{K} f (-\Delta)^{-2s}f \dd \mu.
	\end{align}
	The proof now follows from \cref{convergence}.
\end{proof}

\section{Simulations}
As a part of this thesis we have created and simulated a discrete fractional Gaussian field on the Vicsek set $V_m$. The precise field we have constructed was the one given in \cref{OneToSimulate}. 
These simulation have had a varying rate of success, as we have been able to simulate $X_s^{m}(x) = \sum\lambda_j \phi_j$ for a reasonable set of $s$-values, but we have hit a bottleneck in the fact that we are not able to simulate $X_s^{m}$ for $m \geq 7$. The reason for this bottleneck lies in \cref{discreteMatrix}. 

As stated in the remark for any function $f: V_m \to \ree$ we can define a matrix $L$ such that
\begin{equation}
	L_mF = \begin{pmatrix}
		\Delta_m f(x_1)\\
		\vdots \\
		\Delta_{m}f(x_{N_m})
	\end{pmatrix}, \quad \text{where } F = \begin{pmatrix}
	f(x_1)\\
	\vdots \\
	f(x_{N_m})
	\end{pmatrix}.
\end{equation}
It is precisely from this matrix $L_m$ we compute the eigenvalues and eigenfunctions of $-\Delta_m$. One shall now note that the amount of points in $V_m$ is $5^{m}\cd 4 +1$ so as $m$ grows this number explodes (see \ref{TableOfValues}), and there is not any smart way to construct the matrix $L_m$\footnote{That we have found while researching \textcolor{red}{Formulation of this?}}. Thus the method of constructing the matrix $L_m$ was simply by constructing the Vicsek set $V_m$ and manually checking each point for neighbours.
\begin{table}[!ht]\label{TableOfValues}
	\centering
	\begin{tabular}{|c|c|}
		\hline
		$m$  & $\# V_m$ \\ \hline
		$1$  & $21$\\ \hline
		$2$  & $101$\\ \hline
		$3$  & $501$\\ \hline
		$4$  & $2501$\\ \hline
		$5$  & $12501$\\ \hline
		$6$  & $62501$ \\ \hline
		$7$  & $312501$\\ \hline
		$10$ & $39 062 501$\\ \hline
		$20$ & $3.81\cdot10^{14}$\\ \hline                 
	\end{tabular}
	\caption{Points in $V_m$}
\end{table}

As we see from the table, the values explode in size as $m$ grows. It is therefore really slow having to check each point. Further the matrix $L_m$ that we construct becomes very large, very fast. Since it needs to hold every permutation of points, that means $L_m$ is a $5^{m}\cd 4 + 1 \times 5^{m}\cd 4 + 1$ matrix. It does have the property that it is a sparse matrix, since there is at most five entries in every column. But for $m= 9$ this matrix still fills well over 1 terrabyte in memory.

Having explained the limitations of the simulation, we will quickly delve into the code that generates the simulations.


\subsection{Construction of the Vicsek set for simulation}
While the construction of the Vicsek set is not terrible with regards to speed, it is also not good. The way it was done is by first generating the first "plus" $V_0$. Then going from $V_0$ to $V_1$ we find the endpoints (Tips of the plus) of $V_0$ and then just add $V_0$ to it. So it becomes  

\begin{lstlisting}[language=python]
def VicsekSet(n, endpoints, points):
	if n <= 0:
		return
	
	for endPoint in endpoints:
		for point in points:
			newPoints.add(endPoint +2*point)
	
	for endPoint in endPoints:
		newEndpoints.add(3*endpoint)
	
	return newPoints + VicsekSet(n-1, newEndPoints, newPoints)
\end{lstlisting}

\begin{lstlisting}[language=python]
def findPointsForLineCollection(self, endPoints: list = None, points: list = None, n: int = None, init_length: int = None) -> list:
	
	
	# Setup if none is given as argument
	if n == None:
	n = self.n
	if init_length == None:
	init_length = 3
	if points == None:
	points = [[np.array([-init_length, 0]), np.array([init_length, 0])], [np.array([0, init_length]), np.array([0,-init_length])]]
	
	# This is to check if the recursion given is smaller than zero
	if n <= 0:
	# This checks if we are in the first iteration or if we have just reached recursion depth zero
	if (([[np.array([-init_length, 0]), np.array([init_length, 0])]] and [[np.array([0, init_length]), np.array([0,-init_length])]]) in points):
	return []
	else:
	return [[np.array([-init_length, 0]), np.array([init_length, 0])], [np.array([0, init_length]), np.array([0,-init_length])]]
	else:
	collection = []
	
	# Setup if none is given as argument
	if endPoints == None:
	endPoints = [np.array([-init_length, 0]), np.array([init_length, 0]), 
								np.array([0, init_length]), np.array([0,-init_length])]
	
	# Finds all the points in the current iteration
	for x in endPoints:
	collection.extend([[y + x*2 for y in lists] for lists in points])
	
	# Find the new endpoints
	newEndPoints = [x + 2*x for x in endPoints]
	
	return collection + self.findPointsForLineCollection(newEndPoints, collection + points, n-1, init_length)
\end{lstlisting}

\begin{enumerate}
	\item The rækkefølge af $x_i \in V_m$ er den samme hver gang på grund af sets.
\end{enumerate}

I have some pretty pictures. Want to se the fields $X_s^{m}(x) = \sum\lambda_j \phi_j$ to $\widetilde{X}_s$?

